\cotefile{80}
\scanpage{55}
Soient $a,b,c$ les sommets de $K_J$, $\alpha,\beta,\gamma$
les sommets ``opposés'' dans $K_{J'}=K_I$, intersections
de $D_a$ avec la ps. dr. joignant $b$ \unsure{:|à} $c$ etc ;
\unsure{soit|soient} $\lambda_a,\lambda_b,\lambda_c$ les arêtes de $K_I$ \unsure{contenant|conten$^{\mathrm{t}}$}
$\alpha,\beta,\gamma$, alors $\{\lambda_a,\lambda_b,\lambda_c\}$ est comp. avec
l'orientation.

En résumé

\uline{Théorème} Le foncteur canonique
du groupoïde des arrang$^{\mathrm{mts}}$ de
\unsure{ps.droites|pseudodroites} sans bigone, vers \inserted{celui} des arrangements
1 dimensionnels de pseudo-droites (sans
bigone) est pleinement fidèle.

Les conditions $(A1)$ à $(A6)$ donnent
en principe une voie pour déterminer
l'image essentielle -- mais il faut \unsure{d'|et|dt}
doit arriver à simplifier les conditions
en question. En principe, il faudrait
encore rajouter la relation d'EP (23) --
mais je soupçonne que cette dernière
est automatique.
% archive sheet number: 54
